% !TeX root = main.tex
\chapter*{ПЕРЕЛІК УМОВНИХ ПОЗНАЧЕНЬ, СИМВОЛІВ, ОДИНИЦЬ, СКОРОЧЕНЬ І ТЕРМІНІВ}
\addcontentsline{toc}{chapter}{ПЕРЕЛІК УМОВНИХ ПОЗНАЧЕНЬ, СИМВОЛІВ, ОДИНИЦЬ, СКОРОЧЕНЬ І ТЕРМІНІВ}

\begingroup
\small
\begin{longtable}{>{\bfseries}p{0.20\textwidth}p{0.72\textwidth}}
AUC & площа під кривою залежності повноти від частки хибнопозитивних рішень;\\
Acc (Accuracy) & частка правильних рішень;\\
DCT & дискретне косинусне перетворення;\\
DWT & дискретне хвилькове перетворення;\\
EXIF & формат службових даних зображення;\\
FN & кількість хибнонегативних рішень;\\
FNR & частка хибнонегативних рішень;\\
FP & кількість хибнопозитивних рішень;\\
FPR & частка хибнопозитивних рішень;\\
IPTC & стандарт полів службових даних мультимедійних об’єктів;\\
JPEG & формат стиснення та зберігання зображень;\\
LSB & найменш значущий біт;\\
MP4 & мультимедійний контейнер на основі ISO Base Media File Format;\\
PNG & формат растрових зображень без втрат;\\
Prec (Precision) & точність позитивного виявлення;\\
Rec (Recall) & повнота виявлення;\\
ROC & крива залежності повноти від частки хибнопозитивних рішень;\\
F1 & гармонійне середнє точності позитивного виявлення та повноти;\\
TN & кількість правильно визначених об’єктів норми;\\
TP & кількість правильно визначених аномальних об’єктів;\\
ViT & модель зорового перетворювача;\\
WebP & формат растрових зображень зі стисненням із втратами або без втрат;\\
XMP & розширювана платформа службових даних.\\
\end{longtable}
\endgroup

\clearpage
